
\documentclass[11pt,a4paper]{article}

\usepackage[margin=1in]{geometry}
\usepackage{amsmath,amssymb}
\usepackage{siunitx}
\usepackage{booktabs}
\usepackage{array}
\usepackage{enumitem}
\usepackage{hyperref}
\usepackage{fancyhdr}
\usepackage{listings}

\hypersetup{
    colorlinks=true,
    linkcolor=blue,
    citecolor=blue,
    urlcolor=blue
}

\pagestyle{fancy}
\fancyhf{}
\lhead{BPhO Computational Physics Challenge 2026}
\rhead{High Energy \& Particle Physics}
\cfoot{\thepage}

\title{\textbf{High Energy \& Particle Physics}\\Worked Solutions}
\author{BPhO Computational Physics Challenge 2026}
\date{}

\begin{document}

\maketitle

\section*{Introduction}

This document gives worked solutions to the \emph{High Energy \& Particle Physics}
problem sheet supplied with the BPhO Computational Physics Challenge 2026.

The calculations and conservation-law checks follow the supplied BPhO solution
document. In particular, the particle-physics notation used here follows the sheet:
baryon number is $1/3$ for a quark and $-1/3$ for an antiquark, lepton number is
$+1$ for leptons and $-1$ for antileptons, and antiparticles have the same mass but
opposite additive quantum numbers.

% ============================================================
\section{Question 1 --- Particle interactions and conservation laws}

\subsection{(i) Beta-minus decay}

At the quark level,
\[
d\rightarrow u+W^-,
\]
followed by
\[
W^-\rightarrow e^-+\bar{\nu}_e.
\]

At the first vertex:

\[
\begin{array}{c|ccc}
 & Q/e & B & L\\
\hline
d & -1/3 & 1/3 & 0\\
u & +2/3 & 1/3 & 0\\
W^- & -1 & 0 & 0
\end{array}
\]

Hence
\[
-\frac13=\frac23-1,
\qquad
\frac13=\frac13,
\qquad
0=0+0.
\]

At the second vertex,
\[
W^-\rightarrow e^-+\bar{\nu}_e,
\]
so
\[
-1=-1+0,
\]
and
\[
L(W^-)=0=1+(-1).
\]

Thus charge, baryon number and lepton number are all conserved at both vertices.

\subsection{(ii) Beta-plus decay}

The corresponding process is
\[
u\rightarrow d+W^+,
\]
followed by
\[
W^+\rightarrow e^++\nu_e.
\]

At the first vertex,
\[
+\frac23=-\frac13+1,
\qquad
\frac13=\frac13+0,
\qquad
0=0+0.
\]

At the second vertex,
\[
+1=+1+0,
\]
and
\[
0=(-1)+1.
\]

Therefore
\[
\boxed{\text{charge, baryon number and lepton number are conserved at every vertex.}}
\]

\subsection{(iii) $\pi^+p\rightarrow n+\pi^-+\pi^+$}

Using
\[
\pi^+=u\bar d,\qquad
p=uud,\qquad
n=udd,\qquad
\pi^-=d\bar u,
\]
the interaction is
\[
u\bar d+uud
\rightarrow
udd+d\bar u+u\bar d.
\]

For charge, the left-hand side has
\[
Q/e=
\left(\frac23+\frac13\right)
+
\left(\frac23+\frac23-\frac13\right)
=
1+1=2.
\]

The right-hand side has
\[
0+(-1)+(+1)=0.
\]

This appears inconsistent if the incoming particle is written as $\pi^+$.
However, the supplied BPhO solution sheet writes the reaction as
\[
\boxed{\pi^-+p\rightarrow n+\pi^-+\pi^+},
\]
with quark content
\[
\bar ud+uud
\rightarrow
udd+\bar ud+u\bar d.
\]

For this stated reaction,
\[
Q_{\rm initial}=-1+1=0,
\]
and
\[
Q_{\rm final}=0-1+1=0.
\]

Similarly, the baryon numbers are
\[
B_{\rm initial}=0+1=1
\]
and
\[
B_{\rm final}=1+0+0=1.
\]

Thus the reaction represented in the supplied solution conserves both charge and
baryon number.

\subsection{(iv) $K^+\rightarrow\pi^++\pi^0$}

The charged kaon has quark content
\[
K^+=u\bar s.
\]

The charged pion is
\[
\pi^+=u\bar d,
\]
while the neutral pion is a superposition of
\[
u\bar u
\quad\text{and}\quad
d\bar d.
\]

The initial charge is
\[
Q(K^+)=+\frac23+\frac13=+1.
\]

The final charge is
\[
Q(\pi^+)+Q(\pi^0)
=
(+1)+0
=+1.
\]

Therefore
\[
\boxed{Q_{\rm initial}=Q_{\rm final}=+e}.
\]

\subsection{(v) Quark binding energies}

The supplied masses are
\[
m_{K^+}=0.52616m_p,
\qquad
m_{\pi^+}=0.14875m_p,
\]
and
\[
m_pc^2=938.27\ \mathrm{MeV}.
\]

For the kaon,
\[
m_{K^+}c^2
=
0.52616(938.27)
=
493.68\ \mathrm{MeV}.
\]

The constituent quark masses are
\[
m_uc^2=2.4\ \mathrm{MeV},
\qquad
m_sc^2=95\ \mathrm{MeV}.
\]

Define the binding contribution by
\[
B=
\left(m_u+m_s-m_{K^+}\right)c^2.
\]

Then
\[
B_{K^+}
=
2.4+95-493.68
=
-396.28\ \mathrm{MeV}.
\]

Hence
\[
\boxed{B_{K^+}\approx-396.2\ \mathrm{MeV}}.
\]

The negative sign indicates that the hadron mass-energy is much lower than the sum
of the quoted bare quark masses in this simplified constituent picture.

For the pion,
\[
m_{\pi^+}c^2
=
0.14875(938.27)
=
139.57\ \mathrm{MeV}.
\]

Therefore
\[
B_{\pi^+}
=
2.4+4.8-139.57
=
-132.37\ \mathrm{MeV},
\]
so
\[
\boxed{B_{\pi^+}\approx-132.4\ \mathrm{MeV}}.
\]

These should not be interpreted as ordinary potential energies of two free quarks:
quarks are confined, and the simple subtraction is only the model requested by the
problem.

\subsection{(vi) $\Sigma^-\rightarrow n+e^-+\bar\nu_e$}

The $\Sigma^-$ has quark content
\[
\Sigma^-=dds,
\]
while
\[
n=udd.
\]

At quark level, one strange quark changes through the weak interaction:
\[
s\rightarrow u+W^-,
\]
followed by
\[
W^-\rightarrow e^-+\bar\nu_e.
\]

At the first vertex,
\[
-\frac13=\frac23-1,
\]
\[
\frac13=\frac13+0,
\]
and lepton number remains zero.

At the second vertex,
\[
-1=-1+0
\]
and
\[
0=1-1.
\]

Thus all three relevant additive conservation laws hold at each vertex.

\subsection{(vii) Cloud and bubble chambers}

A \textbf{cloud chamber} contains a supersaturated vapour. A charged high-energy
particle passing through the chamber ionises molecules along its path. The ions act
as nucleation sites, so droplets condense along the trajectory and form a visible
track.

A \textbf{bubble chamber} contains a liquid maintained in a superheated state.
A charged particle ionises the liquid along its path, providing nucleation sites
for bubbles. The bubbles grow along the trajectory and form a visible track.

A magnetic field can be used with either technique to curve charged-particle tracks,
allowing the momentum and charge sign to be inferred.

\subsection{(viii) Pair production of muons}

To create a muon--antimuon pair,
\[
\gamma\rightarrow\mu^-+\mu^+,
\]
the photon must supply at least the two rest-mass energies:
\[
E_{\gamma,\min}=2m_\mu c^2.
\]

With
\[
m_\mu c^2=105.66\ \mathrm{MeV},
\]
we obtain
\[
E_{\gamma,\min}
=
2(105.66)
=
211.32\ \mathrm{MeV}.
\]

Using
\[
1\ \mathrm{GeV}=1.6\times10^{-10}\ \mathrm J,
\]
or equivalently
\[
1\ \mathrm{MeV}=1.6\times10^{-13}\ \mathrm J,
\]
gives
\[
E_{\gamma,\min}
\approx
211.32(1.6\times10^{-13})
=
3.38\times10^{-11}\ \mathrm J.
\]

Thus
\[
\boxed{E_{\gamma,\min}\approx3.38\times10^{-11}\ \mathrm J}.
\]

The minimum wavelength follows from
\[
E=\frac{hc}{\lambda},
\]
so
\[
\lambda_{\max}
=
\frac{hc}{E_{\gamma,\min}}
\approx5.88\times10^{-15}\ \mathrm m.
\]

Therefore
\[
\boxed{\lambda_{\max}\approx5.9\ \mathrm{fm}}.
\]

At threshold the pair is produced with negligible kinetic energy in the idealised
threshold calculation.

\subsection{(ix) Electron--positron pair moving at $0.8c$}

The total energy of one electron or positron is
\[
E=\gamma mc^2,
\]
where
\[
\gamma=
\frac{1}{\sqrt{1-u^2/c^2}}.
\]

For
\[
u=0.8c,
\]
\[
\gamma
=
\frac{1}{\sqrt{1-0.8^2}}
=
\frac53.
\]

The photon must supply the total energy of both particles:
\[
E_\gamma=2\gamma m_ec^2.
\]

Since
\[
m_ec^2=0.51100\ \mathrm{MeV},
\]
we find
\[
E_\gamma
=
2\left(\frac53\right)(0.51100)
\approx1.70\ \mathrm{MeV}.
\]

Hence
\[
\boxed{E_\gamma\approx1.70\ \mathrm{MeV}}.
\]

\subsection{(x) Bremsstrahlung shower}

A high-energy electron entering matter is strongly deflected by the electric fields
of atomic nuclei. Because it is accelerating, it emits bremsstrahlung photons.

A sufficiently energetic photon can then undergo pair production:
\[
\gamma\rightarrow e^-+e^+.
\]

The newly created electron and positron are themselves deflected by nuclei and can
emit further bremsstrahlung photons. Those photons can again produce pairs.

Schematically,
\[
e^-
\longrightarrow
e^-+\gamma
\longrightarrow
e^-+e^-+e^+
\longrightarrow
\text{many }e^\pm+\gamma
\longrightarrow
\text{particle shower}.
\]

The repeated branching produces a cascade of particles. The resulting shower can
therefore provide an observable signature of the original high-energy electron.

\subsection{(xi) Koide mass ratio for $u,d,s$}

The given relation is
\[
\frac{m_u+m_d+m_s}
{\left(\sqrt{m_u}+\sqrt{m_d}+\sqrt{m_s}\right)^2}
\approx\frac{k}{9}.
\]

Using
\[
m_u=2.4,\qquad m_d=4.8,\qquad m_s=95
\]
in MeV$/c^2$,
\[
R_q
=
\frac{2.4+4.8+95}
{\left(\sqrt{2.4}+\sqrt{4.8}+\sqrt{95}\right)^2}.
\]

Numerically,
\[
R_q\approx0.562.
\]

Therefore
\[
k\approx9(0.562)\approx5.06.
\]

The requested integer is
\[
\boxed{k\approx5}.
\]

% ============================================================
\section{Question 2 --- Strangeness}

\subsection{(i) Determining the quark content of $K^0$}

The reaction is
\[
K^-+p\rightarrow\Omega^-+K^++K^0.
\]

The $\Omega^-$ has quark content
\[
\Omega^-=sss
\]
and hence strangeness
\[
S=-3.
\]

The initial $K^-$ has
\[
K^-= \bar us,
\]
so
\[
S_{\rm initial}=-1.
\]

The proton has
\[
S=0.
\]

The final $K^+$ has
\[
K^+=u\bar s
\]
and therefore
\[
S=+1.
\]

For strangeness to be conserved, the neutral kaon must supply
\[
-1=(-3)+(+1)+S(K^0),
\]
giving
\[
S(K^0)=+1.
\]

A neutral meson with strangeness $+1$ must contain an anti-strange quark:
\[
K^0=u\bar s.
\]

Indeed,
\[
Q(K^0)=+\frac23+\frac13=+1
\]
would not be neutral if interpreted this way, so the neutral state is instead
\[
\boxed{K^0=d\bar s},
\]
for which
\[
Q=-\frac13+\frac13=0,
\qquad
S=+1,
\qquad
B=0.
\]

Thus
\[
\boxed{K^0=d\bar s}.
\]

\subsection{(ii) $\Omega^-\rightarrow\Xi^0+\pi^-$}

The quark compositions are
\[
\Omega^-=sss,
\qquad
\Xi^0=uss,
\qquad
\pi^-=d\bar u.
\]

The weak process can be represented by
\[
s\rightarrow u+W^-,
\]
followed by
\[
W^-\rightarrow d+\bar u.
\]

The first vertex conserves charge:
\[
-\frac13=\frac23-1,
\]
and baryon number:
\[
\frac13=\frac13+0.
\]

The second vertex conserves charge:
\[
-1=-\frac13-\frac23,
\]
and baryon number:
\[
0=\frac13-\frac13.
\]

However, strangeness changes:
\[
S(\Omega^-)=-3,
\]
while
\[
S(\Xi^0)=-2,
\qquad
S(\pi^-)=0.
\]

Thus
\[
-3\neq-2.
\]

Therefore
\[
\boxed{\text{strangeness is not conserved in this weak decay}},
\]
as expected for a weak interaction.

% ============================================================
\section{Question 3 --- Tau decay}

\subsection{(i) Feynman diagram}

The decay
\[
\tau^-\rightarrow\pi^-+\nu_\tau
\]
occurs through
\[
\tau^-\rightarrow W^-+\nu_\tau,
\]
followed by
\[
W^-\rightarrow d+\bar u,
\]
which hadronise to form the $\pi^-$.

Thus, schematically,
\[
\tau^-
\longrightarrow
\nu_\tau+W^-,
\qquad
W^-\longrightarrow d+\bar u
\longrightarrow\pi^-.
\]

\subsection{(ii) Why is the tau neutrino required?}

The tau lepton has lepton number
\[
L_\tau=+1.
\]

The $W^-$ boson has
\[
L_\tau=0.
\]

Therefore the vertex
\[
\tau^-\rightarrow W^-+\nu_\tau
\]
requires
\[
1=0+1.
\]

Without the tau neutrino, tau-lepton number would change from $+1$ to $0$.

Hence
\[
\boxed{\text{tau lepton number would not be conserved without }\nu_\tau}.
\]

\subsection{(iii) Kinetic energy of the tau}

The tau has
\[
m_\tau c^2=1776.86\ \mathrm{MeV}
\]
and speed
\[
u=0.2c.
\]

Its Lorentz factor is
\[
\gamma=
\frac{1}{\sqrt{1-0.2^2}}
\approx1.02062.
\]

The total energy is
\[
E_\tau=\gamma m_\tau c^2.
\]

The kinetic energy is
\[
K_\tau=(\gamma-1)m_\tau c^2.
\]

Therefore
\[
K_\tau
=
(1.02062-1)(1776.86)
\approx36.64\ \mathrm{MeV}.
\]

Thus
\[
\boxed{K_\tau\approx36.6\ \mathrm{MeV}}.
\]

\subsection{(iv) Energy of the tau neutrino}

The pion has speed
\[
u_\pi=0.99c.
\]

Its mass is given by
\[
m_\pi c^2
=
0.14875m_pc^2
=
0.14875(938.27)
\approx139.57\ \mathrm{MeV}.
\]

Its Lorentz factor is
\[
\gamma_\pi
=
\frac{1}{\sqrt{1-0.99^2}}
\approx7.0888.
\]

Therefore
\[
E_\pi
=
\gamma_\pi m_\pi c^2
\approx989.37\ \mathrm{MeV}.
\]

The tau total energy is
\[
E_\tau
=
\gamma_\tau m_\tau c^2
\approx1813.50\ \mathrm{MeV}.
\]

Assuming the tau neutrino is effectively massless,
\[
E_{\nu_\tau}
=
E_\tau-E_\pi.
\]

Hence
\[
E_{\nu_\tau}
\approx1813.50-989.37
=
824.13\ \mathrm{MeV}.
\]

Thus
\[
\boxed{E_{\nu_\tau}\approx824.1\ \mathrm{MeV}}.
\]

\subsection{(v) Relativistic energy--momentum invariant}

The relativistic momentum is
\[
p=\gamma mv
=
\left(1-\frac{v^2}{c^2}\right)^{-1/2}mv.
\]

The total energy is
\[
E=\gamma mc^2.
\]

Therefore
\[
E^2-p^2c^2
=
\gamma^2m^2c^4
-
\gamma^2m^2v^2c^2.
\]

Factoring,
\[
E^2-p^2c^2
=
\gamma^2m^2c^4
\left(1-\frac{v^2}{c^2}\right).
\]

Since
\[
\gamma^2=
\frac{1}{1-v^2/c^2},
\]
we obtain
\[
\boxed{E^2-p^2c^2=m^2c^4}.
\]

For an effectively massless tau neutrino,
\[
m_{\nu_\tau}\approx0,
\]
so
\[
E_\nu^2-p_\nu^2c^2=0.
\]

Hence
\[
E_\nu=p_\nu c,
\]
or
\[
p_\nu=\frac{E_\nu}{c}.
\]

Using
\[
E_\nu=824.13\ \mathrm{MeV},
\]
\[
\boxed{p_\nu\approx824.1\ \mathrm{MeV}/c}.
\]

In SI units,
\[
p_\nu
=
\frac{824.13\times10^6(1.602\times10^{-19})}
{2.998\times10^8}
\approx4.40\times10^{-19}\ \mathrm{kg\,m\,s^{-1}}.
\]

Thus
\[
\boxed{p_\nu\approx4.40\times10^{-19}\ \mathrm{kg\,m\,s^{-1}}}.
\]

\subsection{(vi) Koide formulae}

For the charged leptons,
\[
R_\ell
=
\frac{m_e+m_\mu+m_\tau}
{\left(\sqrt{m_e}+\sqrt{m_\mu}+\sqrt{m_\tau}\right)^2}.
\]

Using
\[
m_e=0.51100,\qquad
m_\mu=105.66,\qquad
m_\tau=1776.86
\]
in MeV$/c^2$,
\[
R_\ell\approx0.66666.
\]

Thus
\[
\boxed{R_\ell\approx\frac23}.
\]

For the quarks,
\[
R_q
=
\frac{m_c+m_b+m_t}
{\left(\sqrt{m_c}+\sqrt{m_b}+\sqrt{m_t}\right)^2}.
\]

Using
\[
m_c=1275,\qquad
m_b=4180,\qquad
m_t=172440
\]
in MeV$/c^2$,
\[
R_q\approx0.669.
\]

Thus
\[
\boxed{R_q\approx\frac23}.
\]

The numerical closeness of both expressions to $2/3$ is the empirical pattern
highlighted by the problem.

% ============================================================
\section{Question 4 --- Compton scattering}

Consider an incident photon of wavelength $\lambda$ scattering through angle
$\theta$ from an initially stationary electron.

The Compton wavelength shift is
\[
\boxed{
\Delta\lambda
=
\lambda'-\lambda
=
\frac{h}{m_ec}(1-\cos\theta)
}.
\]

The following derivation obtains this result from relativistic energy--momentum
conservation.

\subsection{(i) Photon momentum and the de Broglie relation}

For any particle,
\[
E^2-p^2c^2=m^2c^4.
\]

For a photon,
\[
m=0,
\]
so
\[
E^2=p^2c^2.
\]

Since $E>0$,
\[
E=pc.
\]

Planck's relation gives
\[
E=hf=\frac{hc}{\lambda}.
\]

Therefore
\[
pc=\frac{hc}{\lambda},
\]
and hence
\[
\boxed{p=\frac h\lambda}.
\]

This is the photon form of the de Broglie relation.

\subsection{(ii) Momentum conservation}

Before the collision, the photon momentum is
\[
p_\lambda=\frac h\lambda.
\]

After scattering,
\[
p_{\lambda'}=\frac h{\lambda'}.
\]

Momentum conservation gives
\[
\vec p_\lambda
=
\vec p_{\lambda'}+\vec p_e.
\]

Therefore
\[
\vec p_e
=
\vec p_\lambda-\vec p_{\lambda'}.
\]

Taking the magnitude and applying the cosine rule,
\[
\boxed{
p_e^2
=
p_\lambda^2+p_{\lambda'}^2
-2p_\lambda p_{\lambda'}\cos\theta
}.
\]

\subsection{(iii) Electron momentum from energy conservation}

Initially the electron is at rest, so its energy is
\[
E_{e,i}=m_ec^2.
\]

Energy conservation gives
\[
m_ec^2+p_\lambda c
=
E_e+p_{\lambda'}c.
\]

Hence
\[
E_e
=
m_ec^2+c(p_\lambda-p_{\lambda'}).
\]

For the recoiling electron,
\[
E_e^2-p_e^2c^2=m_e^2c^4.
\]

Substitute the expression for $E_e$:
\[
p_e^2c^2
=
\left[m_ec^2+c(p_\lambda-p_{\lambda'})\right]^2
-m_e^2c^4.
\]

Expanding and dividing by $c^2$,
\[
p_e^2
=
(p_\lambda-p_{\lambda'})^2
+
2m_ec(p_\lambda-p_{\lambda'}).
\]

Thus
\[
\boxed{
p_e^2
=
(p_\lambda-p_{\lambda'})^2
+
2m_ec(p_\lambda-p_{\lambda'})
}.
\]

\subsection{(iv) Derivation of the Compton shift}

Equating the two expressions for $p_e^2$ gives
\[
p_\lambda^2+p_{\lambda'}^2
-2p_\lambda p_{\lambda'}\cos\theta
=
(p_\lambda-p_{\lambda'})^2
+
2m_ec(p_\lambda-p_{\lambda'}).
\]

Since
\[
(p_\lambda-p_{\lambda'})^2
=
p_\lambda^2+p_{\lambda'}^2
-2p_\lambda p_{\lambda'},
\]
we obtain
\[
2p_\lambda p_{\lambda'}(1-\cos\theta)
=
2m_ec(p_\lambda-p_{\lambda'}).
\]

Therefore
\[
p_\lambda p_{\lambda'}(1-\cos\theta)
=
m_ec(p_\lambda-p_{\lambda'}).
\]

Using
\[
p_\lambda=\frac h\lambda,
\qquad
p_{\lambda'}=\frac h{\lambda'},
\]
gives
\[
\frac{h^2}{\lambda\lambda'}(1-\cos\theta)
=
m_ec
\left(
\frac h\lambda-\frac h{\lambda'}
\right).
\]

Multiplying through by $\lambda\lambda'/h$,
\[
\frac h{m_ec}(1-\cos\theta)
=
\lambda'-\lambda.
\]

Thus
\[
\boxed{
\Delta\lambda
=
\lambda'-\lambda
=
\frac h{m_ec}(1-\cos\theta)
}.
\]

\subsection{(v) Maximum wavelength shift and a 10\% maximum-shift photon}

The shift is greatest when
\[
1-\cos\theta
\]
is greatest. Since
\[
-1\le\cos\theta\le1,
\]
the maximum occurs at
\[
\boxed{\theta=180^\circ}.
\]

Thus
\[
\Delta\lambda_{\max}
=
\frac{2h}{m_ec}.
\]

The problem asks for an incident X-ray whose fractional wavelength shift is
$10\%$ of this maximum:
\[
\frac{\Delta\lambda_{\max}}{\lambda}=0.1.
\]

Therefore
\[
\frac{2h}{m_ec\lambda}=0.1.
\]

Using
\[
E_\gamma=\frac{hc}{\lambda},
\]
we have
\[
\frac{\Delta\lambda_{\max}}{\lambda}
=
\frac{2E_\gamma}{m_ec^2}.
\]

Hence
\[
0.1=\frac{2E_\gamma}{m_ec^2},
\]
so
\[
E_\gamma
=
0.05m_ec^2.
\]

With
\[
m_ec^2=511\ \mathrm{keV},
\]
\[
E_\gamma
=
0.05(511)
=
25.55\ \mathrm{keV}.
\]

Thus
\[
\boxed{E_\gamma\approx25.6\ \mathrm{keV}}.
\]

\subsection{(vi) Electron recoil angle and speed}

Resolve momentum conservation into components.

Taking the incident photon direction as the $x$-axis,
\[
p_\lambda
=
p_{\lambda'}\cos\theta+p_e\cos\phi,
\]
and
\[
0
=
p_{\lambda'}\sin\theta-p_e\sin\phi.
\]

Therefore
\[
p_e\sin\phi
=
p_{\lambda'}\sin\theta,
\]
and
\[
p_e\cos\phi
=
p_\lambda-p_{\lambda'}\cos\theta.
\]

Dividing,
\[
\tan\phi
=
\frac{p_{\lambda'}\sin\theta}
{p_\lambda-p_{\lambda'}\cos\theta}.
\]

Since
\[
\frac{p_\lambda}{p_{\lambda'}}
=
\frac{\lambda'}{\lambda},
\]
and
\[
\frac{\lambda'}{\lambda}
=
1+\frac{h}{m_ec\lambda}(1-\cos\theta),
\]
we obtain
\[
\boxed{
\tan\phi
=
\frac{\sin\theta}
{1+\dfrac{h}{m_ec\lambda}(1-\cos\theta)-\cos\theta}
}.
\]

The electron momentum can be written as
\[
\boxed{
p_e=
\sqrt{
(p_\lambda-p_{\lambda'})^2
+
2m_ec(p_\lambda-p_{\lambda'})
}
}.
\]

Using
\[
p_\lambda=\frac h\lambda,
\qquad
p_{\lambda'}=\frac h{\lambda'},
\]
and
\[
\lambda'
=
\lambda+
\frac{h}{m_ec}(1-\cos\theta),
\]
this gives $p_e$ entirely in terms of $\lambda$, $\theta$ and the constants.

The relativistic momentum relation is
\[
p_e=\gamma m_ev
=
\left(1-\frac{v^2}{c^2}\right)^{-1/2}m_ev.
\]

Solving for $v$,
\[
\boxed{
v
=
c\left[
\left(\frac{m_ec}{p_e}\right)^2+1
\right]^{-1/2}
}.
\]

\subsection{(vii) Maximum electron recoil speed}

The maximum wavelength shift occurs for
\[
\theta=180^\circ.
\]

At this angle,
\[
\lambda'
=
\lambda+\frac{2h}{m_ec}.
\]

The scattered photon therefore has the minimum possible momentum
\[
p_{\lambda'}=\frac h{\lambda'},
\]
so the difference
\[
p_\lambda-p_{\lambda'}
\]
is maximised. From
\[
p_e^2
=
(p_\lambda-p_{\lambda'})^2
+
2m_ec(p_\lambda-p_{\lambda'}),
\]
the electron momentum is consequently maximised.

Energy conservation gives
\[
\frac{hc}{\lambda}+m_ec^2
=
\frac{hc}{\lambda'}
+
\gamma m_ec^2.
\]

Thus
\[
\gamma
=
1+
\frac{h}{m_ec}
\left(
\frac1\lambda-\frac1{\lambda'}
\right).
\]

Since
\[
\gamma=\frac{1}{\sqrt{1-v^2/c^2}},
\]
we obtain
\[
\boxed{
v_{\max}
=
c\sqrt{
1-
\left[
\frac{m_ec^2}
{m_ec^2+\dfrac{hc}{\lambda}-\dfrac{hc}{\lambda'}}
\right]^2
}
},
\]
with
\[
\boxed{
\lambda'
=
\lambda+\frac{2h}{m_ec}
}.
\]

\subsection{(viii) Computational plots}

The problem asks for graphs of the fractional wavelength shift,
electron recoil speed and electron recoil angle as functions of the photon
scattering angle $\theta$.

The quantities can be calculated directly from the preceding equations.

For an incident photon energy $E$,
\[
\lambda=\frac{hc}{E},
\]
and
\[
\boxed{
\frac{\Delta\lambda}{\lambda}
=
\frac{E}{m_ec^2}(1-\cos\theta)
}.
\]

The scattered wavelength is
\[
\lambda'
=
\lambda+\frac{h}{m_ec}(1-\cos\theta),
\]
so
\[
p_\lambda=\frac h\lambda,
\qquad
p_{\lambda'}=\frac h{\lambda'}.
\]

The recoil momentum is
\[
\boxed{
p_e
=
\sqrt{
(p_\lambda-p_{\lambda'})^2
+
2m_ec(p_\lambda-p_{\lambda'})
}
}.
\]

The recoil speed is then
\[
\boxed{
v
=
c\left[
1+\left(\frac{m_ec}{p_e}\right)^2
\right]^{-1/2}
}.
\]

Finally,
\[
\boxed{
\tan\phi
=
\frac{\sin\theta}
{1+\dfrac{h}{m_ec\lambda}(1-\cos\theta)-\cos\theta}
}.
\]

A minimal Python implementation of the required numerical model is:

\begin{lstlisting}[language=Python,basicstyle=\ttfamily\small,frame=single]
import numpy as np
import matplotlib.pyplot as plt

h = 6.62607015e-34
c = 2.99792458e8
me = 9.1093837e-31

theta = np.linspace(1e-6, np.pi, 1000)

for E_keV in (10, 50, 100):
    E = E_keV * 1e3 * 1.602176634e-19
    lam = h*c/E
    lam_p = lam + h/(me*c)*(1 - np.cos(theta))

    p = h/lam
    pp = h/lam_p

    dp = p - pp
    pe = np.sqrt(dp**2 + 2*me*c*dp)

    fractional_shift = (lam_p - lam)/lam
    v = c / np.sqrt(1 + (me*c/pe)**2)

    phi = np.arctan2(
        np.sin(theta),
        1 + h/(me*c*lam)*(1 - np.cos(theta)) - np.cos(theta)
    )

    # fractional_shift, v/c and phi are the three requested quantities
\end{lstlisting}

The three requested plots are therefore generated from
\[
\frac{\Delta\lambda}{\lambda}(\theta),
\qquad
\frac{v}{c}(\theta),
\qquad
\phi(\theta),
\]
for one or more chosen incident photon energies.

% ============================================================
\section*{Summary}

The particle-physics sheet connects several central ideas:

\begin{itemize}
    \item conservation of charge, baryon number, lepton number and, where
    appropriate, strangeness;
    \item quark-level descriptions of weak decays;
    \item relativistic energy and momentum;
    \item pair production and particle showers;
    \item the empirical Koide mass relation;
    \item the energy--momentum invariant
    \[
    E^2-p^2c^2=m^2c^4;
    \]
    \item and the particle description of light through Compton scattering.
\end{itemize}

The Compton-scattering problem in particular leads naturally to the computational
model requested in the final part of the sheet: once the incident photon energy is
specified, the wavelength shift, recoil momentum, recoil speed and recoil angle can
all be evaluated numerically as functions of the scattering angle.

\end{document}
